\section{The Residual Stream vs.\ Propositions: Superposition and the $D_{ff}/D$ Ratio}
\label{sec:superposition}

The formal BP construction uses $D_{\mathrm{model}} = 8$ with one dedicated dimension per
proposition and one OR gate per belief update. A production transformer has $D_{\mathrm{model}}
= 4096$ and $D_{ff} = 14{,}336$: approximately 3.5 OR gates per residual stream dimension per
layer. This raises two related questions. Why does the ratio $D_{ff}/D$ exceed 1 at all? And
does $D_{ff} \approx 14{,}336$ mean there are 14,336 factor nodes per layer? This section
resolves both.

\subsection*{Hidden Units Are Not Factor Nodes}

In the minimal construction, one FFN hidden unit implements one factor node --- one
conditional relationship between two propositions. Hidden units and factor nodes are in
one-to-one correspondence.

In a production transformer, this correspondence breaks down in two ways:

\begin{enumerate}
\item \textbf{Multiple units per factor.} A single conditional relationship may require
multiple hidden units to represent accurately with non-sigmoid activations. ReLU and SwiGLU
require more units to approximate the exact log-odds combination that sigmoid implements in
one unit. The factor node is still the unit of semantics; the hidden units are the
implementation substrate.

\item \textbf{Superposed factors.} Because the residual stream represents many concepts via
superposition, a single hidden unit may participate in multiple factor relationships
simultaneously. The unit fires for a direction in residual stream space that is a
superposition of several feature directions, and the contribution it writes is a superposition
of several belief updates.
\end{enumerate}

The right level of abstraction is not hidden units but \emph{SAE features}. Each monosemantic
SAE feature corresponds to one factor node in the implicit factor graph. The number of
semantically distinct factor nodes per layer is the number of active SAE features --- typically
much smaller than $D_{ff}$.

$D_{ff}$ is therefore the \emph{implementation budget} for the OR gates, not the factor node
count. It sets an upper bound on representable conditional relationships per layer, but the
actual semantic content is determined by the SAE feature structure.

\subsection*{Three Reasons the Budget Exceeds One}

Three independent factors drive $D_{ff}/D$ above the minimal value of 1:

\textbf{Superposition.} Elhage et al.~\cite{elhage2022superposition} show that residual
streams represent far more features than $D_{\mathrm{model}}$ dimensions via superposition:
features are encoded as nearly-orthogonal directions in a high-dimensional space, with
interference managed by sparsity. A model with $D = 4096$ may track tens of thousands of
distinct concepts simultaneously. If the residual stream tracks $K$ active concepts per token
and each needs an independent belief update per layer, $D_{ff}$ must be at least $K$. Since
$K \gg D$, the OR budget must exceed the residual stream dimension by the superposition
density --- empirically, roughly $3$--$4\times$.

\textbf{$k$-Ary neighbors.} The formal construction handles pairwise factors: each OR gate
takes exactly two inputs. Real knowledge graph nodes have more than two neighbors. The scaling
theorem of~\cite{coppola2026transformers} handles this: any $k$-ary factor graph binarizes
exactly via a chain of $\lceil \log_2 k \rceil$ pairwise updates. A node with 8 neighbors
requires 3 chained pairwise updates rather than 1, contributing a further multiplicative
overhead to the required $D_{ff}$.

\textbf{Activation precision.} Sigmoid implements the log-odds algebra exactly. ReLU and
SwiGLU are compatible but not exact: they preserve non-negativity and smooth aggregation but
require more hidden units to achieve the same approximation quality. The bayes-learner
experiment confirms that gradient descent finds approximately correct BP solutions even with
non-sigmoid activations (val MAE 0.000752), but at the cost of additional hidden units per
factor.

Together these three factors explain why the minimal $D_{ff}/D = 1$ becomes $3$--$4$ in
production models. The ratio is an architectural constant encoding the combined implementation
overhead of superposition density, neighbor fanout, and activation imprecision.

\subsection*{The Dimension Budget and Factor Graph Scale}

A cleaner way to think about the implicit factor graph uses the dimension budget directly.
The residual stream has $D_{\mathrm{model}}$ dimensions. Each active concept occupies a
direction in this space. Empirically, roughly $10 \times D_{\mathrm{model}}$ features can be
superposed with manageable interference~\cite{elhage2022superposition}.

For Llama~3 405B with $D_{\mathrm{model}} = 16{,}384$, this gives roughly 160,000
simultaneously active features per token position. Each feature is a factor graph node. The
factor graph has on the order of $W \times 160{,}000 \approx 1.3\mathrm{B}$ nodes for a
sequence of length 8192. This is the scale of the implicit Bayesian network that one forward
pass of Llama~3 405B is performing inference over.

\subsection*{The OR/AND Asymmetry}

With this clarification, the $D_{ff}/H \approx 416$ ratio for Llama~3 405B has a clean
interpretation. It is not that there are 416 OR gates for every AND gate. It is that the OR
implementation requires 416 units of hidden dimension for every unit of attention head
dimension. The factor graph itself may be much sparser --- each node may have only 2--5
neighbors in practice --- but implementing those sparse connections accurately in the
superposed representation requires the $416\times$ overhead.

The implication for interpretability: finding the factor graph requires working at the SAE
feature level, not the hidden unit level. The factor potentials are encoded in the relationship
between input SAE features (what the $W_1$ rows respond to) and output SAE features (what the
$W_2$ columns write). The $D_{ff}$ hidden units are the implementation substrate; the SAE
features are the semantics.